% Chapter introduction: why the topic arose, who first asked, what it answers.
% STATUS: written 2026-11-05 (all five body sections sourced and gate-clean).
% The introduction reports the five threads of the chapter and cites only the
% sources that carry them (Kadison 1990, Crist 1996, Ayupov--Kudaybergenov
% 2012, Chebotar--Ke--Lee 2004, Alaminos--Bresar--Extremera--Villena 2009,
% Liu--Zhang 2017, Yang--Fang 2019). The classical first sources are
% atlas-pending.
\emph{A check at each point.} The early chapters define a derivation by
one law at every pair of points. This chapter weakens the law to a check
at a single point or a single pair, and asks how much of the derivation
survives. Five threads answer, each read from the papers that carry it.

\hbfigure{08-local-characterizations/local}{A map $\Delta$ is a local derivation if for each $x$ there is a derivation $D_x$ with $\Delta(x)=D_x(x)$; the derivation may change from point to point. Is $\Delta$ itself a derivation?}{fig:local:local}

The first thread is the local derivation of Kadison
\cite[p.~494]{Kadison1990Local}. Kadison started the line in 1990 with
von Neumann algebras, and Section~\ref{sec:local-characterizations:local}
reads his Theorem~A, the Larson--Sourour companion and the CSL theorems of
Crist \cite[p.~83]{Crist1996Local}. On every algebra these papers study,
a local check already gives a derivation.

The second thread strengthens the check to two points
\cite[p.~15]{AyupovKudaybergenov2012}. \v{S}emrl asked in 1997 whether a
map that agrees with a derivation at each pair of points is a derivation;
Section~\ref{sec:local-characterizations:two-local} recalls his theorem
through Ayupov and Kudaybergenov and runs their ladder up to every von
Neumann algebra \cite[p.~2]{AyupovKudaybergenov2015}. A two-point check
again gives a derivation on these algebras.

The third thread checks the law only at zero products
\cite[p.~218]{ChebotarKeLee2004}. Chebotar, Ke and Lee moved the
derivation law to the zero products of a prime ring
\cite[p.~218]{ChebotarKeLee2004}.
Section~\ref{sec:local-characterizations:zero-products} reads their
Theorem~2. It also reads the derivational points of Pan on nest algebras and
the zero product determined algebras of Brice and Huang.

The fourth thread comes from preserver problems
\cite[p.~131]{AlaminosBresar2009}. Alaminos, Bre\v{s}ar, Extremera and
Villena started from maps that send zero products to zero products
\cite[p.~131]{AlaminosBresar2009}.
Section~\ref{sec:local-characterizations:preservers} reads their bilinear
tool and their Theorem~2.11 and Theorem~4.6. The answer: a map derivable
at zero products is a generalized derivation.

The fifth thread changes the law into the Lie law
\cite[p.~209]{LiuZhang2017LocalLie}.
Section~\ref{sec:local-characterizations:local-lie-jordan} reads the local
Lie theorem of Liu and Zhang. It then reads the four-condition pattern of
He, Li, An and Huang and the structure theorem of Yang and Fang
\cite[p.~244]{YangFang2019}. The statement: a linear 2-local Lie
derivation on a finite von Neumann algebra without central summands of
type $I_1$ is a Lie derivation \cite[p.~250]{YangFang2019}.
